\documentclass[12pt]{article}
\usepackage[a4paper,margin=2.5cm]{geometry}
\usepackage[T1]{fontenc}
\usepackage[utf8]{inputenc}
\usepackage{lmodern}
\usepackage{microtype}
\usepackage{needspace}
\usepackage{booktabs}
\usepackage{amsmath}
\usepackage{xcolor}
\usepackage{listings}
\usepackage[hidelinks]{hyperref}
\definecolor{codegray}{RGB}{245,245,245}
\lstset{basicstyle=\ttfamily\small,backgroundcolor=\color{codegray},frame=single,
  breaklines=true,showstringspaces=false,columns=fullflexible,keepspaces=true}
\setlength{\parindent}{0pt}
\setlength{\parskip}{0.55em}

\begin{document}
\pagenumbering{gobble}
\begin{titlepage}
\centering
{\large Ministerul Educa\c{t}iei \c{s}i Cercet\u{a}rii al Republicii Moldova\par}
\vspace{0.45cm}
{\large Universitatea Tehnic\u{a} a Moldovei\par}
\vspace{0.45cm}
{\large Facultatea Calculatoare, Informatic\u{a} \c{s}i Microelectronic\u{a}\par}
\vspace{0.45cm}
{\large Departamentul Ingineria Software \c{s}i Automatic\u{a}\par}
\vspace{3cm}
{\Large\bfseries Report on Laboratory Work No. 13\par}
\vspace{0.55cm}
{\LARGE\bfseries Semaphores in Your Own Kernel\par}
\vspace{0.35cm}
{\large Operating Systems\par}
\vfill
\begin{flushleft}
\textbf{Done by:} st. gr. FAF-241, Ca\c{t} Chiril\\[0.8cm]
\textbf{Checked by:} asist. univ. Rei\c{t}man Patricia
\end{flushleft}
\vfill
{Chi\c{s}in\u{a}u -- 2026\par}
\end{titlepage}
\pagenumbering{roman}
\tableofcontents
\clearpage
\pagenumbering{arabic}

\section{Aim, environment, and deliverables}
In Labs 11 and 12, applications used synchronization objects supplied by
Linux's libraries. This lab implements counting semaphores inside xv6 as
three new system calls. A semaphore combines a protected integer with
wait-channel registration and scheduler sleep/wakeup. All three required
parts are complete in the integrated \texttt{xv6-riscv/} tree.

The tools were Ubuntu 24.04.5 LTS on aarch64, RISC-V GCC 13.3.0, and QEMU
8.2.2. The xv6 guest ran with three CPUs and 128 MiB RAM. I verified that
the starting tree built and booted, then captured the user-side checkpoint
before adding the handlers. The final implementation was rebuilt from clean
sources using xv6's \texttt{-Wall -Werror} compilation settings.

\begin{center}
\begin{tabular}{lll}
\toprule
Call & Number & Meaning \\
\midrule
\texttt{sem\_init(id, value)} & 24 & Set the initial counter \\
\texttt{sem\_down(id)} & 25 & Take a permit, waiting if needed \\
\texttt{sem\_up(id)} & 26 & Return a permit and wake waiters \\
\bottomrule
\end{tabular}
\end{center}
The existing \texttt{sync} (22) and \texttt{trace} (23) calls remain intact.
IDs 0--7 identify entries in one global kernel table shared by all
processes; no per-process semaphore copy is made by fork. The optional
\texttt{sem\_trydown} syscall 27 was not implemented.

The integrated tree is the working submission. Fourteen matching changed
source files, the Makefile included, are copied under \texttt{lab13/code/}.
The patch records changes against the pre-Lab-13 tree. The required
\texttt{notes.txt} contains the four Observe answers. Editable report
sources, the assignment, actual runs, and verification are retained under
\texttt{report\_files/}. The additional \texttt{semcheck}, \texttt{semsleep},
and \texttt{semprobe} user programs are regression/evidence tools, not
additional system calls.

\begin{lstlisting}
cd xv6-riscv
make clean
make -j4 TOOLPREFIX=riscv64-linux-gnu- kernel/kernel fs.img
make qemu TOOLPREFIX=riscv64-linux-gnu- CPUS=3
# Exit QEMU with Ctrl-A followed by X.
\end{lstlisting}

\clearpage
\section{Part 1: syscall numbers and the user side}
\subsection{Following the number into the kernel}
I added numbers 24--26 in \texttt{kernel/syscall.h}, declarations in
\texttt{user/user.h}, and three \texttt{entry(...)} lines in
\texttt{user/usys.pl}. The generated assembly stubs load the call number
into \texttt{a7}, run \texttt{ecall}, and return. The first two arguments
are passed in \texttt{a0} and \texttt{a1}; the return value comes back in
\texttt{a0}.

The Makefile includes \texttt{semtest} and \texttt{mix} in \texttt{UPROGS},
so their binaries are placed on the xv6 disk image. The submitted versions
follow the handout: \texttt{semtest} checks order and ping/pong behavior;
\texttt{mix} writes each character separately and optionally protects a
whole line with semaphore 3.

\subsection{Check 1: reaching unknown calls}
After the user-side changes, before kernel dispatch entries existed, the
first lines of the actual checkpoint were:
\lstinputlisting[basicstyle=\ttfamily\footnotesize]{output/checkpoint1-excerpt.txt}
The child ran before the parent because the unknown down call returned
an error instead of blocking. The later checkpoint output contains more
unknown-call errors and interleaved ping/pong characters. The complete
original capture is retained separately; this excerpt is explicitly
shortened.

This proves the new user stubs reached the kernel with the intended
numbers. It is the required intermediate failure, not the final behavior.
The handout's \texttt{semtest} does not check every syscall return value,
so its final ``OK'' string can appear even in this broken intermediate
state. Final verification therefore checks the entire expected ordering,
absence of unknown calls, and independent return-value regressions rather
than relying on that string alone.

\subsection{User programs and semaphore roles}
For the first test, semaphore 0 starts at zero: the child must wait for
the parent's post. Ping/pong uses semaphore 1 at one and semaphore 2 at
zero, passing a turn permit back and forth. Semaphore 3 starts at one
for the optional line lock in \texttt{mix}. Each program initializes its
IDs before forking or starting the corresponding exchange.

\clearpage
\section{Part 2: the kernel semaphore table}
\subsection{Counter, spinlock, and dispatch}
The new \texttt{kernel/sem.c} contains eight static entries:
\begin{lstlisting}[language=C]
struct sem {
  struct spinlock lock;
  int value;
};
static struct sem sems[NSEM];
\end{lstlisting}
A boot-time \texttt{seminit()} explicitly initializes every spinlock and
zeroes the counts before user processes begin. \texttt{main.c} invokes it
and \texttt{defs.h} declares it. The Makefile links \texttt{sem.o}; the
syscall dispatch table registers all three handlers. The Lab 10 trace name
table is extended too, avoiding missing names when tracing the new calls.

Each handler decodes its integer arguments with \texttt{argint}. IDs outside
0--7 are rejected. Init also rejects a negative start value, then updates
the count under the entry's spinlock. Up increments under the same lock
and calls \texttt{wakeup(s)}. An added overflow guard rejects an increment
at \texttt{0x7fffffff}, preserving the signed counter instead of wrapping.

\subsection{The blocking down operation}
The central wait path is:
\begin{lstlisting}[language=C]
acquire(&s->lock);
while (s->value == 0) {
  sleep_prepare(s);
  release(&s->lock);
  sleep_killable();
  acquire(&s->lock);
  if (killed(myproc())) {
    release(&s->lock);
    return -1;
  }
}
s->value--;
release(&s->lock);
return 0;
\end{lstlisting}
The address of the table entry is the wait channel in both prepare and
wakeup. Registration happens before releasing the counter lock. If up
occurs before the waiter actually sleeps, wakeup clears the registered
channel and the sleep path returns without blocking. If the process is
already sleeping, it becomes runnable. This closes the ordinary
check/register/sleep missed-wakeup window.

The submitted \texttt{sleep\_killable()} is a small robustness extension
explained on the next page. It uses this tree's two-stage interface, not
the book's incompatible \texttt{sleep(chan, lock)} call.

\clearpage
\section{Blocking, cancellation, and the required result}
\subsection{A narrowly scoped cancellation guard}
Code review found a separate kill-timing gap: a kill can arrive after
wait registration but before the process changes to \texttt{SLEEPING}.
The existing kill path wakes an already sleeping process but does not
cancel a pending registration in that interval.

I factored the existing sleep body into a private
\texttt{sleep\_common(interruptible)} helper. Ordinary \texttt{sleep()}
passes zero, preserving the original behavior for unrelated kernel waits.
Semaphore waits use \texttt{sleep\_killable()}, which passes one and checks
\texttt{p->killed} under \texttt{p->lock} immediately before entering
\texttt{SLEEPING}. If killed, it clears the pending channel and returns;
down then detects cancellation and returns an error. The same process
lock serializes this decision with kill. This was an inspection-driven
robustness change, not a claim that a failing pre-fix run was captured.

The registration, semaphore-lock release, actual sleep, and counter-lock
reacquisition remain the same two-phase design required by the handout.
A blocked process does not busy-spin waiting for the counter to change.
The spinlock is held only for the short state checks and updates, never
while switching away to the scheduler.

\subsection{Check 2: semtest works}
All three final invocations, within the same final boot, produced the
required sequence:
\lstinputlisting{output/semtest-1.txt}
The parent's ten-tick delay does not let the child run past down: it
remains blocked until the parent posts. Two turn semaphores then enforce
three strict ping-before-pong pairs. The independent focused tests also
validate syscall return values and child exit statuses.

\clearpage
\section{Part 3: a semaphore as the line lock}
\subsection{Without the lock}
Both processes write five lines, one character per \texttt{write} syscall.
The following is the first actual unlocked run, not the sample text:
\lstinputlisting[basicstyle=\ttfamily\footnotesize]{output/mix-1.txt}
Both unlocked runs showed mixed-content lines. Each run still contained
150 uppercase A characters, 150 lowercase b characters, and ten newlines.
The exact interleaving and apparent line lengths depend on scheduling.
A 30-letter line plus its newline costs \textbf{31 write syscalls}; the
whole sequence is not one atomic console operation.

\subsection{With the lock}
\lstinputlisting[basicstyle=\ttfamily\footnotesize]{output/mix-lock-1.txt}
The down before the character loop and up after it make one complete
line, including its newline, a critical section. All three locked runs
produced exactly five complete A lines and five complete b lines. Other
valid runs may alternate whole lines; neither a fixed process order nor
FIFO fairness is promised by this semaphore.

\clearpage
\section{Handout questions and answers}
\begin{enumerate}
\Needspace{9\baselineskip}
\item \textbf{Question:} Without the lock, why do the letters mix? How many system calls does one line cost?
\par\textbf{Answer:} Each character is a separate write, so the two
processes' writes can interleave on the shared console. Thirty letters
and one newline cost 31 write calls. A loop of separate calls is not
implicitly a critical section.

\Needspace{9\baselineskip}
\item \textbf{Question:} With the lock every line is whole. Which two lines of mix.c make one line a critical section?
\par\textbf{Answer:} Use \texttt{sem\_down(3)} before the loop and
\texttt{sem\_up(3)} after it. The initial value is one, so a successful
down reserves the single permit for the whole loop, including the newline.

\Needspace{10\baselineskip}
\item \textbf{Question:} In sem\_down the test is while (s->value == 0), not if. Why can a process wake up and still find 0?
\par\textbf{Answer:} Up wakes every process registered on the entry's
address, but one post adds only one permit. Another process may acquire
the spinlock first and consume it. A wakeup does not reserve a token;
the waiter must recheck the counter under the lock and sleep again if
it remains zero. A killed waiter is separately detected and cancelled.

\Needspace{12\baselineskip}
\item \textbf{Question:} Sem\_down releases the spinlock before sleep. What would go wrong if it slept holding the spinlock?
\par\textbf{Answer:} Up needs that lock to increment and wake, so keeping
it while blocked prevents the operation needed for progress. xv6 also
requires \texttt{sched()} to hold only the process lock, so sleeping with
the semaphore spinlock would violate the invariant and can panic with
\texttt{sched locks}. Register first, release the semaphore lock, sleep,
and reacquire before retesting.
\end{enumerate}

\section{Scope and limitations}
These are simple global teaching semaphores, not named POSIX objects with
ownership or access control. Initialize/reset IDs before concurrent use;
concurrent reinitialization of a live semaphore is outside the demonstrated
contract. There is no fairness guarantee or automatic token recovery if
a process exits while holding a permit. The optional nonblocking syscall
27 was not added. The cancellation and overflow guards do not change
those limits.

\clearpage
\section{Additional verification and conclusion}
\subsection{Focused semantic tests}
The checked \texttt{semcheck} regressions cover all eight IDs; invalid
IDs -1 and 8; negative counts; preservation of the maximum count after
rejected overflow; six children released by six posts; 1000 ping/pong
rounds; and twenty immediate/already-blocked waiter kills. Child statuses
are checked too. Every group passed, ending with
\texttt{semcheck: ALL TESTS PASSED}. These tests do not exhaust all schedules.

\subsection{A sleeping waiter and compatibility with trace}
During \texttt{semsleep}, I sent Ctrl-P to the xv6 console while the child
was blocked on an empty semaphore:
\lstinputlisting[basicstyle=\ttfamily\footnotesize]{output/semsleep.txt}
PID 51 was reported as \texttt{sleep}; posting later woke it and it exited
successfully. This supports the blocking design. It is not a measurement
of total machine CPU use or a CPU-utilization benchmark.

The combined mask for numbers 24--26 is 117440512. The retained trace
probe produced:
\lstinputlisting[basicstyle=\ttfamily\footnotesize]{output/trace-semaphores.txt}
This checks that the existing Lab 10 tracer knows the new names and prints
their zero success results. The old read-only trace still returned
1023, 965, 453, and 0 for the unchanged 2441-byte README.

The \texttt{usertests -q} quick suite ran \textbf{68 tests} and reported
\texttt{ALL TESTS PASSED}; deliberate user traps are not kernel panics.
All \textbf{62 evidence assertions} passed, including required sequences,
character counts, sleeping-state and trace evidence, and byte-for-byte
comparison of fourteen source copies with the tested tree.

A protected count, registration before release, and scheduler sleep turn
semaphore waits into blocking operations rather than busy-waiting. The
loop is essential because a wakeup does not reserve a permit. With one
initial permit, two semaphore calls protect a whole line containing 31
separate writes, connecting the patterns of Labs 11--12 to kernel scheduling.
\end{document}
